\section{Worked example: the board graph and a graph with one edge}\label{sec:example}
% claims used here: C28 C9 m1-5 m1-6 m1-7 m2-13 m2-14 m3-3 m3-4 m3-6 m3-7 m3-8
% Every matrix below is a macro generated by appendices/scripts/generate_paper_numbers.py from the saved script outputs.

This section writes out every matrix for two graphs and applies
\cref{thm:walkcount} and \cref{thm:test} to them. The two graphs are the board graph $B$ (the
matrix on the board, \cref{fig:village}, left) and the graph $D$ with one edge (\cref{fig:oneedge}).
Every matrix and count below was computed by script, and the script outputs were
checked against independent code before the numbers were written
(\cref{app:scripts}).

\subsection{Powers of the board matrix}\label{sec:exB}

\begin{plainwords}
We take the village from the board. It has three houses in a row. The first and
last houses each have one road. The middle house has two roads. We count trips of
two roads, then of three roads. We list every trip, so the reader can see that
each count is right.
\end{plainwords}

\begin{example}[The board matrix and its powers]\label{ex:boardpowers}
The graph $B$ has vertices $n_1,n_2,n_3$ and edges $n_1n_2$ and $n_2n_3$. Its adjacency matrix and powers, copied from the script output, are
\begin{align*}
A&=\MBoardA, & A^2&=\MBoardAsq, & A^3&=\MBoardAcube,\\
A^4&=\MBoardAfour, & A^5&=\MBoardAfive . &&
\end{align*}
The matrix $A$ is symmetric with zero diagonal, as \cref{prop:Aproperties} requires; its row sums $1,2,1$ are the degrees, and the sum of all its entries is $4=2|E|$. Further, $A^3=2A$ and $A^{k+2}=2A^k$ for $k=1,\dots,5$ (by script), and the diagonal of $A^2$ is $(1,2,1)$, the degrees, with sum $4=2|E|$ (\cref{cor:walkcount}(d)). Each power is symmetric (\cref{cor:walkcount}(c)).
\end{example}

The board matrix is the matrix of a graph with two edges. The board drawing also shows an edge $1$--$3$ and an arrowhead; they are not in the matrix (\cref{tab:mismatch}), and this section follows the matrix.

\paragraph{Every walk listed.}
For every entry of $A$, $A^2$ and $A^3$ the script lists the walks and checks that their number equals the entry (labels $1,2,3$ stand for $n_1,n_2,n_3$).
\begin{itemize}
\item Length $1$: $1\text{-}2$ and $2\text{-}3$, and their reverses.
\item Length $2$: from $1$ to $1$: $1\text{-}2\text{-}1$ (one walk). From $1$ to $3$: $1\text{-}2\text{-}3$ (one). From $2$ to $2$: $2\text{-}1\text{-}2$ and $2\text{-}3\text{-}2$ (two). From $3$ to $3$: $3\text{-}2\text{-}3$ (one). All other entries of $A^2$ are $0$.
\item Length $3$: from $1$ to $2$: $1\text{-}2\text{-}1\text{-}2$ and $1\text{-}2\text{-}3\text{-}2$ (two). From $2$ to $3$: $2\text{-}1\text{-}2\text{-}3$ and $2\text{-}3\text{-}2\text{-}3$ (two). The entries $(2,1)$ and $(3,2)$ follow by reversing the walks.
\end{itemize}
The induction step of \cref{thm:walkcount} can be checked on three entries, by the sum over the second-last vertex $n_r$ (script output):
\begin{align*}
(A^2)_{22}&=\sum_r A_{2r}A_{r2}=1\cdot1+0\cdot0+1\cdot1=2,\\
(A^2)_{11}&=\sum_r A_{1r}A_{r1}=0\cdot0+1\cdot1+0\cdot0=1,\\
(A^3)_{12}&=\sum_r (A^2)_{1r}A_{r2}=1\cdot1+0\cdot0+1\cdot1=2.
\end{align*}
In the last line the term $r=1$ is the walk $1\text{-}2\text{-}1$ of length $2$ extended by the step to $n_2$, giving $1\text{-}2\text{-}1\text{-}2$, and the term $r=3$ is the walk $1\text{-}2\text{-}3$ extended by the step to $n_2$, giving $1\text{-}2\text{-}3\text{-}2$. The walk $1\text{-}2\text{-}1$, counted by $(A^2)_{11}$, is neither a path nor a trail: this is the counterexample of \cref{prop:warning}.

\paragraph{Walks, trails and paths by length.}
\Cref{tab:walkcounts} counts all walks, trails, paths and closed walks of each length in $B$, over all pairs of start and end vertices, by enumeration. The totals in the walk column equal the sums of all entries of $A^k$: $4,6,8,12$ for $k=1,2,3,4$ (script). From $k=3$ on there is no trail (it would need three distinct edges, and $B$ has two) and no path (it would need four distinct vertices, and $B$ has three). The walk numbers keep growing; this is the warning of \cref{prop:warning} in numbers.

\begin{table}[ht]
\centering
\caption{All walks, trails, paths and closed walks of length $k$ in the board graph $B$, counted over all start and end vertices (by script). Length $0$ counts the three single vertices.}
\label{tab:walkcounts}
\begin{tabular}{@{}ccccc@{}}
\toprule
$k$ & walks & trails & paths & closed walks \\
\midrule
\TabWalkCountsBody
\bottomrule
\end{tabular}
\end{table}

\subsection{The test on the board graph (connected)}\label{sec:testB}

\begin{plainwords}
Take three houses in a row. A road joins the first and the second house. Another
road joins the second and the third. The first and the third share no road. Can
you still go from one to the other? Yes, through the middle house. We expect the
check to say connected.
\end{plainwords}

\begin{example}[Matrix test on $B$, $N=3$, $N-1=2$]\label{ex:testB}
Every product was computed by script. With $N-1=2$ the sum is $S=A+A^2$, and the second form is $(I+A)^2$:
\begin{align*}
A&=\MBoardA, & A^2&=\MBoardAsq,\\
S&=\MBoardS, & (I+A)^2&=\MBoardIAsq .
\end{align*}
$S$ has no zero entry off the diagonal (in fact none anywhere), so by \cref{thm:test} the graph is connected. $(I+A)^2$ has no zero entry, so by \cref{prop:iplusa} the graph is connected. The two answers agree, and breadth-first search from vertex $1$ reaches $\{1,2,3\}$ and agrees too.
The entry $A_{13}$ is $0$, so no road joins $n_1$ and $n_3$, yet $(A^2)_{13}=0\cdot0+1\cdot1+0\cdot0=1$: exactly one walk of length $2$ joins them, namely $1,2,3$. The entry of $(I+A)^2$ at $(1,3)$ is $1=\binom{2}{2}(A^2)_{13}$, since in formula~\eqref{eq:binomial} both $I$ and $A$ have the entry $0$ at $(1,3)$. As a check, $(I+A)^2=I+2A+A^2$ holds entry by entry (script). \kindtag{derived here}
\end{example}

By \cref{cor:components}(b), the smallest $k$ with $(A^k)_{ij}>0$ is the distance: $d(1,2)=d(2,3)=1$ from $A$, and $d(1,3)=2$ from $A^2$, in agreement with \cref{ex:distance}.

\subsection{The test on the graph with one edge (disconnected)}\label{sec:testD}

\begin{plainwords}
Now take three houses again. Only one road exists, between the first and second
house. The third house stands alone. Nobody can walk to it. We expect the check to
say not connected. The check should also show that the third house is a piece by
itself.
\end{plainwords}

\begin{example}[Matrix test on $D$, $N=3$, $N-1=2$]\label{ex:testD}
The graph $D$ has the single edge $n_1n_2$. From the script,
\begin{align*}
A&=\MOneA, & A^2&=\MOneAsq,\\
S&=\MOneS, & (I+A)^2&=\MOneIAsq .
\end{align*}
The off-diagonal zeros of $S$ are at $(1,3)$, $(2,3)$, $(3,1)$ and $(3,2)$, and the zeros of $(I+A)^2$ are at the same four places. So by \cref{thm:test} and \cref{prop:iplusa} the graph is not connected. By \cref{cor:components}(a), vertices $1$ and $2$ lie in one component and vertex $3$ in another: the components are $\{1,2\}$ and $\{3\}$, which agrees with breadth-first search. Higher powers do not help: the entries $(1,3)$ and $(2,3)$ of $A^k$ are $0$ for $k=1,\dots,6$ (script); for instance
\[
A^3=\MOneAcube=A .
\]
In fact no walk reaches vertex $3$ at all. \kindtag{derived here}
\end{example}

\begin{table}[ht]
\centering
\caption{Outcome of the tests on the two graphs. Both forms of the matrix test and breadth-first search give the same answer.}
\label{tab:outcome}
\small
\begin{tabular}{@{}lccc@{}}
\toprule
Graph & Matrix test (\cref{thm:test}) & Form with $I+A$ (\cref{prop:iplusa}) & Components \\
\midrule
$B$, edges $12$, $23$ & connected & connected & $\{1,2,3\}$ \\
$D$, edge $12$ & not connected & not connected & $\{1,2\}$, $\{3\}$ \\
\bottomrule
\end{tabular}
\end{table}

\paragraph{The test on all small graphs.}
As a check of the statements, not as a proof, the script compared the sum test, the form with $I+A$ and the component statement of \cref{cor:components}(a) with breadth-first search on every labelled simple graph with $N=2,\dots,6$ vertices (\cref{tab:exhaustive}). All agree.

\begin{table}[ht]
\centering
\caption{Exhaustive check over all labelled simple graphs with $N$ vertices (by script). ``Sum test agrees'' counts the graphs on which the test of \cref{thm:test} gave the same answer as breadth-first search. The last column counts the graphs on which the same test with only $N-2$ powers gives a wrong answer (\cref{prop:sharp}).}
\label{tab:exhaustive}
\begin{tabular}{@{}rrrrr@{}}
\toprule
$N$ & graphs & connected & sum test agrees & wrong with $N-2$ powers \\
\midrule
\TabExhaustiveBody
\bottomrule
\end{tabular}
\end{table}
